\section{相关工作}
\label{sec:related}

\subsection{基于超声的肝纤维化分级}

超声自动纤维化评估已有多条研究路线。一条路线为 B 型图像输入增加额外物理
测量,如造影增强微血流动态影像~\cite{microflow2023}或弹性成像衍生的硬度
信息~\cite{swe2018staging},以设备要求换取信号质量。另一条路线仅使用常规
B 型图像:借助生成式增广弥补数据有限~\cite{usfibrosisgan2022},采用更高
频率采集以显露更细纹理~\cite{hfus2025},或利用成对的器官切面以挖掘相对
外观差异~\cite{liverspleen2025}。

这些方法几乎都针对病毒性、酒精性或代谢性肝病的少数离散纤维化分期。血吸虫
场景在目标与表型上均不相同。\emph{日本血吸虫}引起的间隔状、网络状纤维化
与其他病因典型的弥漫性实质改变表现不同~\cite{schistoelasto2018,schistopocus2024},
而且流行区筛查项目所用的分级量表比三或四档的分期决策更细。SFibAI 针对该
场景,将分级化为 36 个序数级别上的后验分布并报告其期
望~\cite{xu2026sfibai};本工作继承该表述,并追问解剖结构能为其带来什么。

\subsection{序数与细粒度预测}

将有序的严重程度当作无序类别处理,会丢弃使分级有意义的那份顺序信息,相关
文献相当丰富。一族方法将序数问题改写为一组耦合的二值决策,或从结构上强制
秩一致~\cite{coral2020},或通过条件概率分解实现~\cite{corn2023}。另一族
将连续量离散化为有序区间并预测逐区间概率分布,如序数深度估
计~\cite{dorn2018}。第三族对预测分布施加单峰正则,使概率质量集中于相邻
级别而非散布于整个量表~\cite{unimodalbeta2021}。

在医学影像中,序数表述已被用于疾病严重程度分级,有时与分割联
合~\cite{sonnet2022};面向视网膜分级亦发展出不确定性感知的变
体~\cite{drgraduate2020}。我们的分级头属于第二族:36 级后验分布,其期望即
为报告分数,并结合继承自 SFibAI 的邻域感知软标签项与边界敏感惩罚
项~\cite{xu2026sfibai}。我们并不主张新的序数表述;细粒度量表是研究解剖
先验的任务场景,它在此重要,恰恰因为细微的级别区分正是额外结构上下文可能
有所贡献之处。

\subsection{解剖信息引导与辅助监督学习}

与我们最接近的一类工作,是向执行另一主任务的网络注入解剖知识。若干策略反复
出现。

\textbf{以解剖为条件的注意力。}注意力门控网络无需显式定位监督即可学会突出
显著区域、抑制无关背景~\cite{attentiongated2019}。Anatomy-XNet 进一步用
解剖分割掩码监督注意力,使胸片疾病分类聚焦于肺与心脏区域~\cite{anatomyxnet2022}。
SynAP-Fib 与其意图相同——将表征引向解剖上有意义的证据——但其病灶信号是
粗边界框的并集而非器官分割,且所得注意力以门控残差的形式注入分级,而非用于
对骨干网络自身的空间注意力重新加权。

\textbf{按解剖划分的架构。}AGMB-Transformer 将特征路由到由解剖定义的分支,
用于根管治疗自动评估~\cite{agmbtransformer2021};解剖引导的级联利用视杯
位于视盘之内这一已知包含关系~\cite{opticdisc2020}。HoVer-Trans 将乳腺超声
组织的水平与垂直分层直接编码进 transformer 的注意力结构~\cite{hovertrans2023}。
这些设计将架构绑定于固定的解剖分解。SynAP-Fib 则保持单一共享骨干,把解剖
部位作为\emph{预测的}六类后验来调制分级,因此没有任何分支与特定部位绑定,
推理时也无需部位元数据。

\textbf{带辅助头的多任务学习。}乳腺超声中的分割--分类联合学习已有研究,
联合学习改善了两个任务的表现~\cite{zhou2020mtlabus,shamtl2021};胎儿脑
影像中亦有姿态估计与组织分割的联合应用~\cite{petsnet2023}。弱监督定位以
低得多的标注代价实现类似效果,从图像级监督中导出病灶证据~\cite{covidweak2020,weakbreast3d2019}。
在这些设计中,辅助分支通常与主任务共享梯度,其收益与干扰因此纠缠:辅助损失
直接塑造共享表征。SynAP-Fib 刻意用双向梯度隔离边界分开两条通路,使每个分支
贡献的是其\emph{输出}作为上下文,而非其梯度作为正则。

\textbf{定位。}就单个要素而言,没有一项是新的:解剖监督注意力、部位条件化
预测、辅助多任务头、弱框监督与门控特征融合均有先例。本文考察的这些紧密
相关研究中,多数在其有利配置中只使用单一辅助信号。我们报告的 $2\times2$
比较回答的是另一个问题:不同种类的解剖先验能否复合;它给出的答案,是任何
单先验研究都无法揭示的。
